% =====================================================================================
%  Chapter 4 — Control Architectures
%  Thesis: Comparison of Control Strategies for Interaction-Force Aware Multirotor Teams
%
%  DRAFT, 23 August 2026 (writing-track item W.8). Source: docs/18_Strategy_Descriptions.md.
%  REVISED 2026-10-14 (thesis-revision run, Phase 4). Markers: % NEW, % CHANGED, % DECISION,
%  % VERIFY, % REMOVED. Plan rows P4-xx refer to docs/thesis_revision/stage3_revision_plan.md.
%
%  THE RULE THIS CHAPTER MUST NOT BREAK.
%  Each strategy is presented in two separated parts: what WE implement (verifiable against
%  firmware_app/src/lib.rs and residual_nn.rs) and the REFERENCE METHOD it follows (cited,
%  with deviations stated). These are never merged into one sentence. Our implementation is
%  not the paper's method, and writing as though it were would misrepresent both -- and
%  would undermine RQ1, whose validity rests on each strategy being a defensible
%  representative of its family rather than a reproduction of a specific paper.
%
%  Implementation status is stated honestly per strategy. Several are not implemented.
%  Do not describe an unimplemented strategy as though it exists.
%
%  Citation baseline: 2026-08-27 (see docs/thesis/CITATION_BASELINE.md).
%  After that date: only NEW cites or EDITED prior-work sentences need re-checking.
%
%  Numbering (current): Strategy 0 = uncompensated geometric baseline; 1 = pure INDI;
%  2 = geometric + learned residual; 3 = feedback linearisation + flatness-preserving residual.
%  No other strategy is compared (supervisor decision); further methods appear only as related
%  work in Chapter 2. (CHANGED P4-22: the earlier note on a seven-strategy numbering is obsolete.)
% =====================================================================================

\chapter{Control Architectures}
\label{ch:control}

Chapter~\ref{ch:model} defined the residual wrench and the two ways of obtaining it. This
chapter describes the control strategies compared in this thesis, and in each case the point
at which knowledge of the interaction enters the control law.

Four controllers: an uncompensated baseline (Strategy~0), flown in every scenario as the
reference condition (``uncompensated'' means that it does not use the residual), and the three strategies compared against it (Strategies~1--3).
Table~\ref{tab:ctrl-summary} lists them, and the rest of the chapter takes them in order.

% DECISION D-05 (H9, lab status): the "Status" column and the phrase "at the time of writing" in the
%   caption report project progress ("flying", "implemented, unflown"). Options: (i) delete the
%   column; (ii) replace by a design statement; (iii) neutral wording "implemented in the flight
%   code: yes / yes / yes / no" (recommended). Text left unchanged.
% CHANGED (P4-16): footnote a on the rotor-speed column (training labels need rotor speeds).
\begin{table}[t]
  \centering
  \caption{The strategies compared, and the implementation status of each at the time of writing.
    ``Knows'' refers to the residual: whether the controller has a measurement of it, a prediction
    of it, both, or neither. ``Enters at'' is the quantity the residual is subtracted from, and
    ``Rotor speed'' says whether the strategy needs rotor-speed sensing at run time.}
  \label{tab:ctrl-summary}
  \small
  \begin{tabular}{@{}clllcc@{}}
    \toprule
    \# & Strategy & Knows & Enters at & Rotor speed & Status \\
    \midrule
    0 & Geometric baseline           & nothing    & ---                        & no  & flying \\
    1 & Pure INDI                    & measured   & $a_{\mathrm{cmd}}$, torque & yes & flying \\
    % CHANGED (P4-21, B4-2): footnote a, training labels need rotor speeds (R07 Sec. 5.4.1).
    2 & Geometric + NN               & predicted  & $a_{\mathrm{cmd}}$         & no\textsuperscript{a}  & implemented, unflown \\
    3 & FBL + NN                     & predicted  & linearising law            & no\textsuperscript{a}  & not implemented \\
    \bottomrule
  \end{tabular}

  \vspace{0.4em}
  \begin{minipage}{0.9\linewidth}\footnotesize
    \textsuperscript{a}~Not at run time. Producing the training labels requires the measured
    residual \eqref{eq:model-ares}, and therefore rotor-speed sensing on the vehicles that collect
    the data, as for the learned method of Cobo-Briesewitz et al.~\cite{cobobriesewitz2026lindi}.
  \end{minipage}
\end{table}

The chapter treats each strategy as its own controller. What makes the comparison fair is not a
claimed shared implementation, but the experimental conditions held constant for every one of them — the same
airframe, trajectories, formations, separations, residual definition and gains frozen per strategy
after tuning, stated precisely in Section~\ref{sec:ctrl-fairness}.

Each strategy is described in two parts, kept deliberately separate: what is implemented here, and
the published method it follows. They follow published \emph{families}, and the deviations from
each source paper are listed in that strategy's section — none is a bit-identical replica.
Strategies whose author code has not yet arrived stay unimplemented rather than guessed. Each family
is represented by one implementation of our own, so results are statements about these
implementations, not about the families.

% -------------------------------------------------------------------------------------
\section{The position and attitude loop}
\label{sec:ctrl-common}

Strategies~0, 1 and~2 (the geometric baseline, pure INDI and geometric+NN) share one geometric
position loop and an attitude law for the commanded thrust direction. Strategy~1 replaces the
attitude law's inner term with an incremental one. The shared equations are stated once in this
section. Strategy~3 uses feedback linearisation instead; it has its own section.

\subsection{Position loop}

The commanded acceleration is
%
\begin{equation}
  a_{\mathrm{cmd}}
  \;=\; \underbrace{a_{\mathrm{ff}}}_{\text{trajectory}}
      + K_p e_p + K_v e_v + K_i \!\int\! e_p \,\mathrm{d}t
      + g e_3
      \;\underbrace{-\; a_{\mathrm{res}}}_{\text{measured}}
      \;\underbrace{-\; \hat{a}_{\mathrm{res}}}_{\text{predicted}},
  \label{eq:ctrl-ades}
\end{equation}
% DECISION D-02 (integral terms): the flight code does not contain the vector integral K_i * int(e_p)
%   in the control law (ENABLE_POSITION_INTEGRAL = false, lib.rs:525; term zeroed in
%   controller_step, L1865-1869). It contains a vertical-only integral k_{i,z} * int(e_z) e_3, enabled at
%   compile time (ENABLE_Z_INTEGRAL = true, lib.rs:535; accumulate_z_integral L1548-1569 with
%   conditional anti-windup and a 1.5 m s limit; term added at L1880), whose gain is a run-time
%   parameter (pos_gains.ki_z): the study configuration sets 16 for the geometric controller
%   (Strategies 0 and 2) and 0 for our INDI (Strategy 1) (crazyflies.yaml:237, study/ours.patch:22-23).
%   Options: (a) rewrite to the code (recommended); (b) keep the general form and state which terms
%   are zero; (c) leave. Equation left unchanged.
%   CLOSING PASS (D-11): the values above for Strategy 1 are those of the 'ours' block (controller_step INDI,
%   ctrl_mode 3). The frozen S1 block of crazyswarm2/crazyflie/config/study_strategies.json selects a
%   different INDI controller (stabilizer.controller 9, controller_omar_indi.c), which subtracts the
%   residual (L260, L268) with matched filters and runs a vertical integral (Kpos_Iz 1.5). Which one
%   Strategy 1 is, is decision D-11.
% DECISION D-01 (residual sign): the code adds res_sign * a_res and res_sign * a_nn
%   (lib.rs:1930-1931); default res_sign = +1 (traj_iface.c:442); study: +1 for Strategy 1, -1 for
%   Strategy 2. The two "minus" signs above hold for Strategy 2 only as configured. Equation left
%   unchanged.
%   CLOSING PASS (D-11): the values above for Strategy 1 are those of the 'ours' block (controller_step INDI,
%   ctrl_mode 3). The frozen S1 block of crazyswarm2/crazyflie/config/study_strategies.json selects a
%   different INDI controller (stabilizer.controller 9, controller_omar_indi.c), which subtracts the
%   residual (L260, L268) with matched filters and runs a vertical integral (Kpos_Iz 1.5). Which one
%   Strategy 1 is, is decision D-11.
%
with $e_p = p_{\mathrm{des}} - p$ and $e_v = v_{\mathrm{des}} - v$, and gains diagonal but distinct
in the horizontal and vertical axes. Here $a_{\mathrm{cmd}}$ includes gravity compensation and the two residual terms. The
trajectory-demanded part $a_{\mathrm{ff}} + K_p e_p + K_v e_v + K_i \!\int\! e_p \,\mathrm{d}t$ is the
$a_{\mathrm{des}}$ of Chapter~\ref{ch:model}. The commanded thrust vector is $m\,a_{\mathrm{cmd}}$; its
magnitude projected onto the body $z$-axis gives the collective thrust, and its direction, together
with the desired yaw, gives the desired attitude $R_{\mathrm{des}}$ through the usual flatness
construction. In words: the thrust vector fixes the direction of the body $z$-axis, and the desired
yaw fixes the rotation about that axis.
% NEW (P4-04): the construction, and the conditioning of the velocity error, from the code.
Concretely, with $\psi_{\mathrm{des}}$ the desired yaw and
$x_c = [\cos\psi_{\mathrm{des}},\, \sin\psi_{\mathrm{des}},\, 0]^\top$ the desired heading, the
columns of $R_{\mathrm{des}}$ are
\begin{equation*}
  z_{\mathrm{des}} = \frac{a_{\mathrm{cmd}}}{\lVert a_{\mathrm{cmd}} \rVert}, \qquad
  y_{\mathrm{des}} = \frac{z_{\mathrm{des}} \times x_c}{\lVert z_{\mathrm{des}} \times x_c \rVert}, \qquad
  x_{\mathrm{des}} = y_{\mathrm{des}} \times z_{\mathrm{des}},
\end{equation*}
the construction of Lee et al.\ for the position-and-heading mode~\cite{lee2010geometricse3}, and the
collective thrust is $f = m\,a_{\mathrm{cmd}}^\top R e_3$, clipped at zero. Before the velocity
error enters $K_v e_v$, its horizontal components are limited to \SI{0.6}{\metre\per\second} in
magnitude and passed through a second-order low-pass filter at \SI{10}{\hertz}; both measures
reduce the effect of a single corrupted motion-capture sample, and the vertical component is left
untouched.\footnote{\texttt{firmware\_app/src/lib.rs}, \texttt{desired\_rot},
\texttt{condition\_vel\_error} and \texttt{controller\_step}.}
% CHANGED: code, lib.rs:desired_rot (L1463-1473), condition_vel_error (L1504-1521) with
%   ENABLE_POSE_GLITCH_GUARD / ENABLE_VEL_FB_FILTER (L564-574), controller_step thrust (L1949-1950).

The last two terms of \eqref{eq:ctrl-ades} are the whole subject of this thesis, and both are
subtracted, for the reason derived in Section~\ref{sec:model-sign}. \emph{Which residual
information is present is what distinguishes Strategies 0, 1 and 2 from one another}, and nothing
else about the position loop changes between them: Strategy~0 has neither term, Strategy~1 the
measured one, and Strategy~2 the predicted one. Strategy~3 injects its prediction elsewhere, in
the feedback-linearising law of Section~\ref{sec:ctrl-s3}.
% DECISION D-01 / D-02: "both are subtracted" and "nothing else about the position loop changes"
%   are not true of the study configuration as it stands: the residual sign and the vertical
%   integral gain differ between Strategy 1 and Strategies 0/2 (see the comments under the equation).

\subsection{Attitude loop}

The attitude error is the standard geometric one,
$e_R = \tfrac{1}{2}\left(R_{\mathrm{des}}^\top R - R^\top R_{\mathrm{des}}\right)^\vee$, with
$e_\omega = \omega - \omega_{\mathrm{des}}$, where $\omega_{\mathrm{des}}$ is the desired body angular
velocity that follows from the desired attitude trajectory. In the geometric configuration the torque is
%
\begin{equation}
  \tau = -K_R^{\mathrm{geo}} e_R - K_\omega^{\mathrm{geo}} e_\omega + \omega \times J\omega
         - k_{i,\mathrm{att}} \textstyle\int e_R \, \mathrm{d}t,
  \label{eq:ctrl-tau-geo}
\end{equation}
% DECISION D-02: the attitude integral is disabled in the flight code (ENABLE_ATTITUDE_INTEGRAL =
%   false, lib.rs:540; k_i,att = 0 in controller_step, L2225). Equation left unchanged.
%
% CHANGED (P4-07): units of the gains.
the third term compensating the gyroscopic coupling. Here $K_R^{\mathrm{geo}}$ and
$K_\omega^{\mathrm{geo}}$ are gains from attitude error and rate error to torque, in
\si{\newton\metre} per unit of $e_R$ and \si{\newton\metre\second} per unit of $e_\omega$; $k_{i,\mathrm{att}}$ is the gain of an
integral of the attitude error. This is a reduced form of the law
in~\cite{lee2010geometrictracking}: it drops that work's angular-feedforward bracket and adds an
integral term the source does not have --- see \eqref{eq:ctrl-tau-lee} and the discussion in
Section~\ref{sec:ctrl-s0}. The symbols of that work map to ours as $k_R \to K_R^{\mathrm{geo}}$,
$k_\Omega \to K_\omega^{\mathrm{geo}}$, $\Omega \to \omega$ and $\Omega_d \to \omega_{\mathrm{des}}$. The INDI configurations replace
\eqref{eq:ctrl-tau-geo} with an incremental law, described in Section~\ref{sec:ctrl-s1}.
% CHANGED: code, traj_iface.c g_indi_kr_geo/kw_geo comments (L266-280, torque units).

\subsection{Runtime selection}
\label{sec:ctrl-runtime}

% CHANGED (P4-08): what the configuration selects, stated completely.
% DECISION D-11: this paragraph and Sec. ctrl-s1 describe Strategy 1 as the INDI of controller_step
%   (both INDI bits on). The frozen S1 block of study_strategies.json selects controller 9 instead.
%   Text left unchanged; see author_decisions.md D-11.
The control law is selected by a parameter at runtime rather than by recompiling: one bit enables
the measured residual in the position loop, another selects the incremental attitude law, and a
third parameter enables the learned prediction. Strategy~0 has all three off. Strategy~1 has both
INDI bits on. Strategy~2 has only the prediction on. All are pushed to the vehicle at connection time
from the configuration, as one frozen parameter set per strategy, together with each strategy's own gains, the sign with which the
residual terms enter (Section~\ref{sec:model-sign}) and the gain of the vertical integral term. The
prediction itself is computed whenever a complete weight set is loaded, whatever the third
parameter says; the parameter decides only whether it enters \eqref{eq:ctrl-ades}.
% CHANGED: code, traj_iface.c PARAM_GROUP indi_gains / pos_gains / rnn (L203-214, L522-582);
%   lib.rs:controller_step (L1752-1756, L1837, L1851-1855, L2175).
% CHANGED (V3): the study keeps one configuration file (crazyswarm2/crazyflie/config/crazyflies.yaml)
%   and one frozen parameter block per strategy (crazyswarm2/crazyflie/config/study_strategies.json,
%   blocks S0, S1, S2, ...; documentation string: checked against the read-back values before
%   take-off), so "a single configuration file" is replaced by "one frozen parameter set per
%   strategy". See D-11 for which controller the S1 block selects.

This matters more than it may appear. Comparing strategies that require different firmware
introduces a confound that is impossible to rule out afterwards, a flash cycle changes more than
the control law, and strategies flown on different days meet different battery states, different
temperatures and different motion-capture calibrations. Runtime selection allows the strategies to
be interleaved within a single session, which is what makes the comparison a paired one.

% -------------------------------------------------------------------------------------
\section{Strategy 0: uncompensated geometric baseline}
\label{sec:ctrl-s0}

\paragraph{Implemented.} Equations~\eqref{eq:ctrl-ades} and \eqref{eq:ctrl-tau-geo} with both
residual terms zero. No knowledge of the interaction whatsoever.
% NEW (P4-06): the desired-rate error used in the code.
The rate error uses the desired body rate directly, $e_\omega = \omega - \omega_{\mathrm{des}}$,
where $\omega_{\mathrm{des}}$ follows from the trajectory's jerk by differential flatness or is taken
from the setpoint; it is not transported into the current body frame by $R^\top R_{\mathrm{des}}$ as
in the source (below). The two definitions coincide when the attitude error is zero.
% CHANGED: code, lib.rs:controller_step geometric branch (L2219-2235), omega_desired (L1417-1428).

\paragraph{Reference.} Geometric tracking on $SE(3)$~\cite{lee2010geometrictracking}, which tracks
a position trajectory together with a single heading direction, four outputs, not a full
six-degree-of-freedom pose. The attitude error of \eqref{eq:ctrl-tau-geo} is that work's
$e_R = \tfrac{1}{2}(R_d^\top R - R^\top R_d)^\vee$, verified against the source rather than
reproduced from memory.
% NEW (P4-05): frame mapping.
Lee et al.\ write the translational dynamics with the third axis pointing \emph{down},
$m\dot v = mge_3 - fRe_3$, and define the position error as $e_x = x - x_d$, the opposite of the
$e_p = p_{\mathrm{des}} - p$ used here~\cite{lee2010geometricse3}. Position-loop terms therefore change
sign when transcribed into this thesis' east-north-up convention, whereas the attitude error and
the torque law, written in body coordinates, carry over unchanged.

\paragraph{The implemented torque law is a reduced form, and is labelled as such.} The law in the
source is
%
\begin{equation}
  M = -k_R e_R - k_\Omega e_\Omega + \Omega \times J\Omega
      - J\!\left( \hat{\Omega}\, R^\top R_d \Omega_d - R^\top R_d \dot{\Omega}_d \right),
  \label{eq:ctrl-tau-lee}
\end{equation}
%
with $e_\Omega = \Omega - R^\top R_d \Omega_d$. Equation~\eqref{eq:ctrl-tau-geo} retains the
proportional, damping and gyroscopic terms but omits the final feedforward bracket, and adds an
integral term on attitude error that the source does not contain. It is therefore a
\emph{reduced} law in the same family, not the torque law of~\cite{lee2010geometrictracking}, and
no stability result proved there transfers to it unexamined. Restoring the feedforward bracket
would be a firmware change rather than a wording change, and is outside the scope of this draft.
% DECISION D-02: "adds an integral term" -- the term exists in the code but is disabled.

\paragraph{What the source proves.} Not global asymptotic tracking. With the attitude error
function $\Psi$, exponential stability is established on $\Psi(0) < 1$, initial attitude errors
below \SI{90}{\degree}, and almost-global exponential attractiveness on $1 \le \Psi < 2$,
below \SI{180}{\degree}. The distinction matters here only in that this thesis claims neither: the
baseline is flown as a reference condition, not as a certified controller.
% NEW: definition of Psi for the reader (R15, eq. 6).
Here $\Psi(R, R_d) = \tfrac{1}{2}\operatorname{tr}(I - R_d^\top R)$, which equals $1 - \cos\theta$
for a rotation by the angle $\theta$ between $R$ and $R_d$, so $\Psi = 1$ at \SI{90}{\degree} and
$\Psi = 2$ at \SI{180}{\degree}~\cite{lee2010geometricse3}.
% VERIFY: the regions are attributed to the CDC paper (no PDF held); they are verified against the
%   held companion arXiv:1003.2005v4 (Propositions 3 and 4 there).

\paragraph{Role.} Strategy~0 is the uncompensated baseline of Table~\ref{tab:ctrl-summary}. It is
flown in every scenario so that ``how much did compensation help'' has a denominator, and so that
the interaction can be separated from the model error discussed in
Section~\ref{sec:model-measure-fail}. It is also the controller under
which all training data is collected, precisely because it does not compensate: a residual observed
under a compensating controller is the residual that survived compensation, which is not what a
predictive model should be fitted to.
% DECISION D-04: the last clause contradicts the measurement property stated in Ch3 Sec. 3.3.1
%   (fourth paragraph) and Shi et al. 2020 Sec. III-C: for a residual formed from measured rotor
%   speeds, compensation changes the commanded thrust but not the residual at a given state. Valid
%   reasons for collecting under Strategy 0 are listed in author_decisions.md. Text left unchanged.

% -------------------------------------------------------------------------------------
\section{Strategy 1: pure INDI}
\label{sec:ctrl-s1}
% DECISION D-11: which INDI implementation this section describes must match the one Strategy 1
%   flies (controller_step INDI, ctrl_mode 3, vs controller 9); see author_decisions.md.

\paragraph{Implemented.} The measured residual of \eqref{eq:model-ares} enters the position loop as
$-a_{\mathrm{res}}$ in \eqref{eq:ctrl-ades}. In the attitude loop the geometric torque law is
replaced by an incremental one. A reference angular acceleration is formed from the geometric
errors,
%
\begin{equation}
  \alpha_{\mathrm{ref}} = \alpha_{\mathrm{des}} - K_R^{\mathrm{INDI}} e_R - K_\omega^{\mathrm{INDI}} e_\omega,
  \label{eq:ctrl-alpharef}
\end{equation}
% DECISION D-01: "enters the position loop as -a_res" holds only for res_sign = -1; the study
%   configuration of Strategy 1 uses +1 (study/ours.patch:13-14).
%   CLOSING PASS (D-11): the values above for Strategy 1 are those of the 'ours' block (controller_step INDI,
%   ctrl_mode 3). The frozen S1 block of crazyswarm2/crazyflie/config/study_strategies.json selects a
%   different INDI controller (stabilizer.controller 9, controller_omar_indi.c), which subtracts the
%   residual (L260, L268) with matched filters and runs a vertical integral (Kpos_Iz 1.5). Which one
%   Strategy 1 is, is decision D-11.
%
where $\alpha_{\mathrm{des}}$ is the feedforward angular acceleration obtained from trajectory snap
by differential flatness. The INDI gains $K_R^{\mathrm{INDI}}$ and $K_\omega^{\mathrm{INDI}}$ map
errors to angular acceleration, in \si{\per\second\squared} and \si{\per\second}, so their units
differ from the geometric gains, which map to torque. The achieved angular acceleration $\alpha_{\mathrm{meas}}$ is obtained by
differentiating the gyro signal and low-pass filtering it, and the torque currently being applied,
$\tau_{\mathrm{cur}}$, is reconstructed from measured rotor speeds. The command is then the
increment
%
\begin{equation}
  \tau = \tau_{\mathrm{cur}} + J\left(\alpha_{\mathrm{ref}} - \alpha_{\mathrm{meas}}\right).
  \label{eq:ctrl-indi-tau}
\end{equation}
%
In words: the controller compares the angular acceleration it wants with the one it actually got, and
changes the torque by the difference. Because \eqref{eq:ctrl-indi-tau} corrects relative to what is actually being achieved, it rejects
any disturbance appearing in $\alpha_{\mathrm{meas}}$ without needing a model of it. The
differentiation and filtering are the source of the delay identified in
Section~\ref{sec:model-measure-fail}; the filter cutoff sets the trade-off between that delay and
the noise admitted from the gyro.

% NEW (P4-10): what the code adds to eq. (ctrl-indi-tau).
In the implementation, $\alpha_{\mathrm{ref}}$ passes through the same second-order Butterworth
filter as $\alpha_{\mathrm{meas}}$ before the difference is taken, so that both terms of the
difference carry the same lag; $\tau_{\mathrm{cur}}$ is obtained by inverting the motor mixer, that
is, from the four rotor thrusts $k_{t,k}\varpi_k^2$, the moment arm and the ratio of rotor drag
torque to thrust; the increment is multiplied by an effectiveness factor that defaults to one;
and the resulting torque can be limited per axis. No gyroscopic term is added: the increment
already contains whatever the measured angular acceleration contains.
% CHANGED: code, lib.rs:controller_step (alpha_ref filter L2075-2080; INDI branch L2175-2218),
%   rpms_to_torque (L1481-1493), clamp_torque (L1651-1664); traj_iface.c g_indi_j_scale (L350),
%   g_indi_clamp_en (L370).
% DECISION D-06: the filter cut-off actually realised (filt_dt_us, filt_prewarp; traj_iface.c:444-478)
%   and whether tau_cur is filtered (filt_tau, default 0, traj_iface.c:342) are configuration facts
%   the author should state for the study configuration.

% NEW (P4-09): where eq. (ctrl-indi-tau) comes from, and its translational counterpart.
\paragraph{Why the increment rejects disturbances.} Write the rotational dynamics
\eqref{eq:model-rot} as $J\dot\omega = \tau + \delta$, where
$\delta = \tau_{\mathrm{res}} - \omega\times J\omega$ collects everything the controller does not
command. At the last sample the vehicle applied $\tau_{\mathrm{cur}}$ and experienced
$\alpha_{\mathrm{meas}}$, so $J\alpha_{\mathrm{meas}} \approx \tau_{\mathrm{cur}} + \delta$, and the
unknown is estimated as $\delta \approx J\alpha_{\mathrm{meas}} - \tau_{\mathrm{cur}}$. Requiring
$J\alpha_{\mathrm{ref}} = \tau + \delta$ and substituting the estimate gives
\eqref{eq:ctrl-indi-tau}. This is the derivation of Tal and Karaman, who state the assumptions it
needs: $\delta$ changes slowly compared with the filter and the sample period, and the measured
quantities are filtered alike~\cite{tal2018indi}. Their analysis also shows that, under these
assumptions, the closed-loop angular acceleration follows its reference with the dynamics of the
motors, and that an error in $J$ affects the transient rather than the final value. A one-axis
example with illustrative numbers: with $J = 1$, a constant unknown torque $\delta = 0.3$ and
$\alpha_{\mathrm{ref}} = 0$, a vehicle that applied $\tau_{\mathrm{cur}} = 0$ measures
$\alpha_{\mathrm{meas}} = 0.3$ and commands $\tau = -0.3$, which cancels $\delta$ at the next step;
with motor lag the cancellation takes as long as the motors need to reach the new command.

The position loop of Strategy~1 is the same idea written without an explicit increment.
Substituting \eqref{eq:model-ares} into \eqref{eq:model-sign} gives
\begin{equation*}
  f R e_3 = m\,a_{\mathrm{des}} + mg e_3 - m\left(a_{\mathrm{meas}} - a_{\mathrm{model}}\right)
          = f_{\mathrm{cur}} R e_3 + m\left(a_{\mathrm{des}} - a_{\mathrm{meas}}\right),
\end{equation*}
where $f_{\mathrm{cur}}$ is the thrust reconstructed from the measured rotor speeds: the new thrust
vector is the current one corrected by the acceleration error, which is the linear-acceleration
increment of Tal and Karaman and of Smeur et al.~\cite{tal2018indi,smeur2017cascadedindi}.
% DECISION D-01: this equivalence holds for the residual SUBTRACTED; with the configured res_sign = +1
%   of Strategy 1 the position term is not this increment.
%   CLOSING PASS (D-11): the values above for Strategy 1 are those of the 'ours' block (controller_step INDI,
%   ctrl_mode 3). The frozen S1 block of crazyswarm2/crazyflie/config/study_strategies.json selects a
%   different INDI controller (stabilizer.controller 9, controller_omar_indi.c), which subtracts the
%   residual (L260, L268) with matched filters and runs a vertical integral (Kpos_Iz 1.5). Which one
%   Strategy 1 is, is decision D-11.

The position and attitude increments are independently selectable, giving four configurations:
geometric, position-INDI only, attitude-INDI only, and full INDI. Strategy~1 refers to the full
configuration, and that is the controller the comparison reports. The intermediate ones are
diagnostic bits used in the two-robot results chapter to attribute an effect to one loop or the
other. They are not additional strategies and will not be reported as Strategy~1.

\paragraph{Reference.} The incremental formulation with flatness-derived
feedforward follows Tal and Karaman~\cite{tal2018indi}; the cascaded arrangement of position and
attitude INDI follows Smeur et al.~\cite{smeur2017cascadedindi}.

% CHANGED (P4-12): the rotor-speed source and the filtering deviations, from the code.
\paragraph{Deviations.} The airframe, gains and filter cutoffs are ours and were identified for
this platform. Tal and Karaman obtain rotor speed from optical encoders~\cite{tal2018indi}; this
implementation uses the rotor speeds that the electronic speed controllers report over the DShot
protocol, or an optical rotor-speed deck, after a plausibility filter (Section~\ref{sec:model-measure}).
Two further deviations concern filtering. Tal and Karaman require identical low-pass filters on the
measured acceleration and on the thrust it is compared with~\cite{tal2018indi}; in the position
loop the implementation filters the accelerometer and, by default, not the rotor-speed thrust. In
the attitude loop they filter the moment reconstructed from rotor speeds~\cite{tal2018indi}, whereas
the implementation filters the reference $\alpha_{\mathrm{ref}}$ and, by default, not
$\tau_{\mathrm{cur}}$. Finally, Tal and Karaman note that the incremental structure removes the
need for integral action in their design~\cite{tal2018indi}. No claim is made that this
implementation reproduces the tracking performance reported in either work, that would require a
like-for-like replication which is not part of this thesis.
% CHANGED: code, traj_iface.c rpm_get_all and g_indi_rpm_source (L508-520, L701-730); lib.rs
%   ENABLE_ACC_PREFILTER (L545) and res_fc default 0 (traj_iface.c:493); filt_tau default 0 (L342).
% DECISION D-02: whether Strategy 1 runs a vertical integral (study configuration: gain 0) should
%   be stated here once D-02 is resolved.
%   CLOSING PASS (D-11): the values above for Strategy 1 are those of the 'ours' block (controller_step INDI,
%   ctrl_mode 3). The frozen S1 block of crazyswarm2/crazyflie/config/study_strategies.json selects a
%   different INDI controller (stabilizer.controller 9, controller_omar_indi.c), which subtracts the
%   residual (L260, L268) with matched filters and runs a vertical integral (Kpos_Iz 1.5). Which one
%   Strategy 1 is, is decision D-11.

\paragraph{Requirements.} Rotor-speed sensing, without which $\tau_{\mathrm{cur}}$ and
$a_{\mathrm{res}}$ are both unavailable and the method degrades silently
(Section~\ref{sec:model-measure-fail}). The method needs no training data; the residual is formed
online from inertial and rotor-speed measurements.

% -------------------------------------------------------------------------------------
\section{Strategy 2: geometric with a learned residual}
\label{sec:ctrl-s2}

% CHANGED (P4-13): evaluation rate, input source and deployment from the code.
\paragraph{Implemented.} The geometric controller of Section~\ref{sec:ctrl-s0} with the deep-sets
prediction $\hat{a}_{\mathrm{res}}$ of \eqref{eq:model-deepsets} subtracted in
\eqref{eq:ctrl-ades}. The network is evaluated onboard at every tenth control step, about
\SI{100}{\hertz}, from the relative states of the neighbours: their positions arrive in the
motion-capture broadcast that the flight stack already uses, so the input requires no additional
communication, and their velocities are differenced onboard (Section~\ref{sec:model-deepsets}).
The prediction is a vertical acceleration; the horizontal residual is not compensated by this
strategy. Weights are trained offline on data collected under Strategy~0 and compiled into the
flight image before flight.
% CHANGED: code, traj_iface.c g_rnn_div (L132); lib.rs:rnn_predict_eval (L839-873), rnn_service
%   (flash build, L922-926); residual_nn.rs:eval (vertical output, L302).
% DECISION D-01: "subtracted" holds for the study configuration of Strategy 2 (res_sign = -1,
%   crazyflies.yaml:227).

The prediction is computed and logged \emph{whether it is enabled or not}. This is what makes the
model's own accuracy measurable independently of its effect on the vehicle: on a flight where the
prediction is not driving the controller, predicted and measured residuals are two independent
quantities and their comparison is meaningful. Once the prediction is driving the controller, the
measured residual is the one that survived compensation and the comparison no longer isolates model
quality. Both flights are therefore needed, and Chapter~\ref{ch:setup} schedules them.
% DECISION D-04: "the measured residual is the one that survived compensation" contradicts Ch3
%   Sec. 3.3.1 (fourth paragraph): the measured residual at a given state does not depend on
%   compensation. What changes when the prediction drives the controller is the state visited (the
%   lower vehicle no longer sags), so prediction-on flights test the model on a different state
%   distribution. Text left unchanged.

\paragraph{Reference.} The approach of predicting the interaction from relative states with a
permutation-invariant network and applying it as feedforward follows the Neural-Swarm
line~\cite{shi2020neuralswarm,shi2022neuralswarm2}.
% NEW (P2-11 cross-reference): the closest published instance of exactly this structure.
The closest published instance of the structure used here, a learned residual added as
feedforward to a geometric $SE(3)$ controller, is a secondary configuration in the hardware
experiments of Hsieh et al., whose main controller is the feedback-linearising one of
Strategy~3~\cite{hsieh2026flatness}.

% CHANGED (P4-15): deviations rewritten for the implemented network.
\paragraph{Deviations.} These are substantial and are stated plainly. The network architecture is
Neural-Swarm2's own (Section~\ref{sec:model-deepsets}), but four things differ from that paper.
First, Neural-Swarm2 applies spectral normalisation, which its abstract associates with stability
and generalisation guarantees and on which the bounded-error assumption of its analysis
rests~\cite{shi2022neuralswarm2}; this implementation is trained without it, and the guarantees
that would follow are correspondingly not claimed. Second, Neural-Swarm2 uses the learned residual
in motion planning as well as control~\cite{shi2022neuralswarm2}, whereas here the trajectories are
planned in advance and the residual is used for control only. Third, Neural-Swarm2's tracking
controller has its own attitude law and an explicit compensation of motor delay~\cite{shi2022neuralswarm2};
here the prediction enters the geometric controller of Section~\ref{sec:ctrl-common} unchanged.
Fourth, the model is trained on this project's two-vehicle data, it predicts only the vertical
component, its output is bounded at \SI{8}{\metre\per\second\squared}, at most three neighbours are
considered, and it is evaluated at about \SI{100}{\hertz} rather than at every control step. This
strategy is a representative of the predictive family, and should be read as such rather than as an
evaluation of Neural-Swarm2.
% REMOVED (P4-15): "The architecture here is a deep-sets model of our own sizing, 987 parameters,
%   rather than a reproduction of any published network. Neural-Swarm2 applies spectral
%   normalisation to obtain stability and generalisation guarantees~\cite{shi2022neuralswarm2};"
%   (wrong for the current code; the guarantee sentence also dropped the abstract/body distinction).

% CHANGED (P4-16): rotor speeds for the labels.
\paragraph{Requirements.} Neighbour relative states; training data. No rotor-speed sensing at run
time; the training labels, being the measured residual, require rotor-speed sensing on the
vehicles that collect the data.

% -------------------------------------------------------------------------------------
\section{Strategy 3: feedback linearisation with a flatness-preserving residual}
\label{sec:ctrl-s3}

% DECISION D-05 (H9): "requested from the authors and had not been received at the time of writing"
%   reports project status. Recommended neutral wording: "Not implemented in the flight code. The
%   reference method is described below." Text left unchanged.
\paragraph{Implementation status.} \textbf{Not implemented.} The controller code was requested from the
authors and had not been received at the time of writing. The strategy does not block the others,
and its absence is a scope limitation rather than a failure of the comparison
(the discussion chapter).

\paragraph{Reference.} Hsieh et al.~\cite{hsieh2026flatness}, building on the theory of Yang et
al.~\cite{yang2025flatnessresiduals}. The distinguishing property is where the prediction enters:
not as a feedforward term added to a desired acceleration, but inside a feedback-linearising law
whose validity depends on the augmented system remaining differentially flat. A generic learned
residual can destroy that flatness; residuals with a lower-triangular structure, a structured
correction inside the state derivative, not a free wrench, preserve it, and with it the original
flat output~\cite{yang2025flatnessresiduals}. That theory is established in simulation on a planar
quadrotor; the hardware demonstration is Hsieh et al.'s.
% NEW: what the structure means concretely in Hsieh et al. (R05 Sec. 3.2).
In Hsieh et al.\ the structure takes a simple form: the learned residual may depend only on the
positions and velocities of the vehicles, not on thrust or attitude, so that thrust and attitude
can still be computed algebraically from the flat outputs, position and yaw, and their
derivatives~\cite{hsieh2026flatness}.

% CHANGED (P4-18): simulation and hardware results separated, as in Ch2 Sec. 2.5.
Its relevance to this comparison is that it is the strongest published counter-argument to treating
``predictive compensation'' as one method. Hsieh et al.'s abstract reports a \SI{31}{\percent}
reduction in average tracking error. On hardware the body reports \SI{24}{\percent} RMSE for
ROM-FBL over nominal FBL and a further \SI{29}{\percent} for learned FBL over ROM-FBL, and an
average RMSE of about \SI{5}{\centi\metre} that the authors describe as matching the reported
performance of the nonlinear model predictive controller of Chee et al.; the direct comparison with
that controller, using the same learned residual, is in simulation, where the flatness-based
controller needs about one twentieth of the computation time~\cite{hsieh2026flatness}. The abstract
adds that tight formation flight was achieved with under \SI{30}{\second} of training data and a
\SI{5}{\milli\second} loop rate; on hardware that loop ran \emph{off-board}, on a laptop streaming
commands over the radio~\cite{hsieh2026flatness}.
Those are the figures a predictive strategy in this class should be measured against. None of them
was produced by a controller running on the vehicle.
% REMOVED (P4-18): "In simulation they match the tracking error published by Chee et al. for
%   nonlinear model predictive control at an order of magnitude less computation, from under 30 s
%   of training data, at a 5 ms loop rate obtained off-board" (conflated the hardware match with
%   published figures and the simulation head-to-head; R05).

\paragraph{Requirements.} Neighbour relative states; training data. No rotor-speed sensing at run
time; as for Strategy~2, the training labels require it.

% -------------------------------------------------------------------------------------
\section{What makes the comparison fair, and what it does not guarantee}
\label{sec:ctrl-fairness}

\paragraph{What is held constant.} The airframe, the trajectories, the formations, the
separations, the residual definition and the estimator. Each strategy is tuned for itself; those
gains are then frozen before the comparison flights. This holds uniformly for every strategy — none
is exempted or held to a looser standard on the tasks or the measurement pipeline.

% NEW (P4-21): what differs between strategies by design.
\paragraph{What differs by design.} Because each strategy is tuned for itself, more than the
interaction information differs between them. The attitude law differs (geometric for Strategies~0
and~2, incremental for Strategy~1); the gain sets differ, and the geometric and incremental gains
are not even in the same units; the filter settings of the incremental law apply only to
Strategy~1; and each strategy's configuration fixes the sign with which residual terms enter and
the gain of the vertical integral term (Section~\ref{sec:ctrl-runtime}). These differences are
listed here so that a difference in outcome is read against all of them, not only against the
residual information.
% DECISION D-01 / D-02: the configured values (Strategy 1: res_sign +1, ki_z 0; Strategies 0/2:
%   res_sign -1 where the network is used, ki_z 16) should be stated here once decided.
%   CLOSING PASS (D-11): the values above for Strategy 1 are those of the 'ours' block (controller_step INDI,
%   ctrl_mode 3). The frozen S1 block of crazyswarm2/crazyflie/config/study_strategies.json selects a
%   different INDI controller (stabilizer.controller 9, controller_omar_indi.c), which subtracts the
%   residual (L260, L268) with matched filters and runs a vertical integral (Kpos_Iz 1.5). Which one
%   Strategy 1 is, is decision D-11.

\paragraph{What this does not guarantee.} Two limitations are worth stating before results are
presented rather than after.

Every strategy follows a published \emph{family}, not a paper reproduced exactly. The deviations
from each source are listed in that strategy's section. Strategies whose author code has not yet
arrived stay unimplemented rather than guessed. Each strategy is its own controller. ``Identical
conditions'' therefore means identical trajectories, formations, hardware and separations. It does
not mean identical control structure or one gain set on every law. A difference in outcome could
reflect the controller structure itself, not only how it uses interaction-force information. That
applies to every strategy.

Geometric tracking, INDI and feedback linearisation are different laws and are tuned separately.
After that tuning, every strategy's gains are frozen for the comparison flights so a later gap is
not mid-campaign retuning. The cost of per-strategy tuning is in scope for RQ4. The discussion
chapter reports where a frozen gain appears to have disadvantaged a strategy.

% -------------------------------------------------------------------------------------
\section{Summary}
\label{sec:ctrl-summary}

Three strategies (1--3: pure INDI, geometric+NN and FBL+NN), plus the uncompensated baseline
(Strategy~0) flown as the reference condition under all of them, each its own controller,
distinguished by what they know about the interaction force, when they know it, and where that
knowledge enters. All three are compared under the same experimental conditions: identical
trajectories, formations, hardware and separations, with each strategy's gains frozen after its own
tuning, not identical control structure.

% DECISION D-05 (H9): the next paragraph reports project progress ("two strategies fly", "one is
%   implemented and unflown"). Recommended neutral wording: "Strategies 0, 1 and 2 are implemented
%   in the flight code; Strategy 3 is not." Text left unchanged.
At the time of writing, two strategies fly (the baseline and pure INDI), one is implemented and
unflown (geometric+NN), and one is not implemented (FBL+NN). Chapter~\ref{ch:setup} describes the experimental
design; the two- and multi-robot results chapters report what was flown.
